% commands.tex — bibliothèque de commandes partagée
% Importer dans le préambule avec % commands.tex — bibliothèque de commandes partagée
% Importer dans le préambule avec % commands.tex — bibliothèque de commandes partagée
% Importer dans le préambule avec % commands.tex — bibliothèque de commandes partagée
% Importer dans le préambule avec \input{commands}
% Requiert : \documentclass déjà déclaré

% ── Police ───────────────────────────────────────────────────────────────────
\usepackage[T1]{fontenc}
\usepackage[utf8]{inputenc}
\usepackage{tgheros}
\renewcommand{\familydefault}{\sfdefault}
\newcommand{\titlefont}{\sffamily}

% ── Mise en page ──────────────────────────────────────────────────────────────
\usepackage[a4paper, top=2cm, bottom=2.2cm, left=2.5cm, right=2.5cm]{geometry}
\usepackage{microtype}
\setlength{\parindent}{0pt}
\setlength{\parskip}{0.6em}
\pagestyle{empty}

% ── Couleurs ─────────────────────────────────────────────────────────────────
\usepackage{xcolor}
\definecolor{C1lo}{HTML}{1E3A8A}  % bleu foncé
\definecolor{C1hi}{HTML}{0D9488}  % turquoise
\definecolor{C2lo}{HTML}{F4801A}  % orange
\definecolor{C2hi}{HTML}{C2185B}  % magenta
\definecolor{C3lo}{HTML}{27AE60}  % vert
\definecolor{C3hi}{HTML}{0D9488}  % turquoise (alias C1hi)
\definecolor{Cdark}{HTML}{1A1A2E}
\definecolor{Clight}{HTML}{F3F4F6}
\definecolor{Ccode}{HTML}{F0F0F0}

\definecolor{Cgrey}{HTML}{868686}
\colorlet{grey}{Cgrey}
\definecolor{Cc1}{HTML}{F0F0F0}
\definecolor{Cc2}{HTML}{F0F0F0}

% ── Mise en page avancée ─────────────────────────────────────────────────────
\usepackage{tabularx}
\usepackage{calc}
\usepackage{wrapfig}

% ── Commandes stat ────────────────────────────────────────────────────────────
%
% Six commandes de mise en page côte-à-côte : grand chiffre + sous-titre
% dans une colonne (gauche ou droite), texte courant dans l'autre colonne.
%
% Principe commun (variantes vtop/hbox)
% ──────────────────────────────────────
%   \hbox to \linewidth{ vtop(stat) vtop(texte) \hss }
%
%   vtop avec \kern0pt : point de référence au sommet (height=0, depth=contenu).
%   \parskip=0pt dans le vtop stat empêche un saut parasite avant le gros chiffre.
%
%   Espacement — variable \statgap
%   ──────────────────────────────
%   \statgap = blanc visible entre le chiffre et le texte courant.
%   Commun aux 6 commandes : changer \statgap recalibre tout.
%   Implémentation vtop/hbox : le blanc est absorbé dans la largeur du vtop-stat
%   (statwidth = natural + \statgap, pas de \hspace explicite entre les vtops).
%   Implémentation wrapfigure : \columnsep = \statgap (gap géré par wrapfig).
%
%   Alignement optique (sommet des caractères)
%   ───────────────────────────────────────────
%   Le vtop-texte est relevé par \raise\dimexpr\statcap-\bodycap\relax pour
%   aligner les sommets de glyphes. Non appliqué aux variantes wrapfigure.
%
% Comportement d'ancrage
% ──────────────────────
%   Les variantes vtop/hbox produisent un \hbox nu. Le caller gère \par et blancs.
%
% Les six commandes
% ─────────────────
%   \statcol    stat GAUCHE, vtop/hbox
%   \statflow   stat GAUCHE, vtop/hbox  (alias fonctionnel de \statcol)
%   \statflowr  stat DROITE, vtop/hbox, raggedleft, chiffre à la marge droite
%   \statcolwf  stat GAUCHE, wrapfigure (float : hors vbox/minipage)
%   \statflowwf stat GAUCHE, wrapfigure (idem, alias de \statcolwf)
%   \statflowrwf stat DROITE, wrapfigure (idem)

\newlength{\statwidth}
\newlength{\statcap}
\newlength{\bodycap}
\newlength{\statgap}
\setlength{\statgap}{1.0cm}

% \statcol{chiffre}{sous-titre}{texte}
\newcommand{\statcol}[3]{%
  \settowidth{\statwidth}{{\fontsize{60}{64}\selectfont\bfseries #1}}%
  \addtolength{\statwidth}{\statgap}%
  \settoheight{\statcap}{{\fontsize{60}{64}\selectfont\bfseries X}}%
  \settoheight{\bodycap}{{\large X}}%
  \noindent\hbox to \linewidth{%
    \vtop{\kern0pt\hsize=\statwidth\relax
      \parskip=0pt\relax
      {\fontsize{60}{64}\selectfont\bfseries #1\par}%
      {\large #2}%
    }%
    \raise\dimexpr\statcap-\bodycap\relax
    \vtop{\kern0pt\hsize=\dimexpr\linewidth-\statwidth\relax
      {\large #3\par}%
    }%
    \hss
  }%
}

% \statflow{chiffre}{sous-titre}{texte}  — stat gauche
\newcommand{\statflow}[3]{%
  \settowidth{\statwidth}{{\fontsize{60}{64}\selectfont\bfseries #1}}%
  \addtolength{\statwidth}{\statgap}%
  \settoheight{\statcap}{{\fontsize{60}{64}\selectfont\bfseries X}}%
  \settoheight{\bodycap}{{\large X}}%
  \noindent\hbox to \linewidth{%
    \vtop{\kern0pt\hsize=\statwidth\relax
      \parskip=0pt\relax
      {\fontsize{60}{64}\selectfont\bfseries #1\par}%
      {\large #2}%
    }%
    \raise\dimexpr\statcap-\bodycap\relax
    \vtop{\kern0pt\hsize=\dimexpr\linewidth-\statwidth\relax
      {\large #3\par}%
    }%
    \hss
  }%
}

% \statflowr{chiffre}{sous-titre}{texte}  — stat droite
\newcommand{\statflowr}[3]{%
  \settowidth{\statwidth}{{\fontsize{60}{64}\selectfont\bfseries #1}}%
  \addtolength{\statwidth}{\statgap}%
  \settoheight{\statcap}{{\fontsize{60}{64}\selectfont\bfseries X}}%
  \settoheight{\bodycap}{{\large X}}%
  \noindent\hbox to \linewidth{%
    \raise\dimexpr\statcap-\bodycap\relax
    \vtop{\kern0pt\hsize=\dimexpr\linewidth-\statwidth\relax
      {\large #3\par}%
    }%
    \vtop{\kern0pt\hsize=\statwidth\relax
      \parskip=0pt\relax
      \raggedleft
      {\fontsize{60}{64}\selectfont\bfseries #1\par}%
      {\large #2\par}%
    }%
    \hss
  }%
}

% ── Variantes wrapfigure ──────────────────────────────────────────────────────
%
% Doublons de statcol/statflow/statflowr utilisant wrapfigure.
% Le texte (#3) s'écoule autour du chiffre ; le contenu qui suit la commande
% dans le document wrappera aussi (comportement natif wrapfigure).
%
% Contrainte : wrapfigure est un float. Il ne fonctionne pas dans un vbox,
% minipage ou toute boîte : le float est alors différé hors de la boîte.
% À réserver au flot normal du document.
%
% Gap : \columnsep mis à \statgap en local.
% \intextsep mis à 0pt pour ne pas creuser de vide au-dessus du chiffre.
% La largeur déclarée à wrapfigure = natural_width du chiffre seul (le gap
% est pris en charge par \columnsep, pas inclus dans la boîte figure).

% \statcolwf{chiffre}{sous-titre}{texte}  — stat gauche, wrapfigure
\newcommand{\statcolwf}[3]{%
  \settowidth{\statwidth}{{\fontsize{60}{64}\selectfont\bfseries #1}}%
  \setlength{\intextsep}{0pt}%
  \setlength{\columnsep}{\statgap}%
  \begin{wrapfigure}{l}{\statwidth}%
    \vtop{\kern0pt\hsize=\statwidth\relax
      \parskip=0pt\relax
      {\fontsize{60}{64}\selectfont\bfseries #1\par}%
      {\large #2}%
    }%
  \end{wrapfigure}%
  {\large #3\par}%
}

% \statflowwf{chiffre}{sous-titre}{texte}  — stat gauche, wrapfigure
\newcommand{\statflowwf}[3]{%
  \settowidth{\statwidth}{{\fontsize{60}{64}\selectfont\bfseries #1}}%
  \setlength{\intextsep}{0pt}%
  \setlength{\columnsep}{\statgap}%
  \begin{wrapfigure}{l}{\statwidth}%
    \vtop{\kern0pt\hsize=\statwidth\relax
      \parskip=0pt\relax
      {\fontsize{60}{64}\selectfont\bfseries #1\par}%
      {\large #2}%
    }%
  \end{wrapfigure}%
  {\large #3\par}%
}

% \statflowrwf{chiffre}{sous-titre}{texte}  — stat droite, wrapfigure
\newcommand{\statflowrwf}[3]{%
  \settowidth{\statwidth}{{\fontsize{60}{64}\selectfont\bfseries #1}}%
  \setlength{\intextsep}{0pt}%
  \setlength{\columnsep}{\statgap}%
  \begin{wrapfigure}{r}{\statwidth}%
    \vtop{\kern0pt\hsize=\statwidth\relax
      \parskip=0pt\relax
      \raggedleft
      {\fontsize{60}{64}\selectfont\bfseries #1\par}%
      {\large #2\par}%
    }%
  \end{wrapfigure}%
  {\large #3\par}%
}

% ── TikZ ─────────────────────────────────────────────────────────────────────
\usepackage{tikz}
\usetikzlibrary{shadings, shadows, fadings}

% \gradtext{couleur-gauche}{couleur-droite}{contenu}
% Texte en dégradé gauche → droite, sur n'importe quel fond.
% Méthode : PDF Tr=7 (clip aux contours des glyphes) + dégradé TikZ.
% Après Q, le tracker CTM de pdflatex est désynchronisé ; on réémet un cm
% compensatoire égal à la largeur de l'image TikZ (convertie en bp).
\makeatletter
\newcommand{\gradtext}[3]{%
  \leavevmode%
  \pdfliteral{q 7 Tr }%
  \rlap{{#3}}%
  \pdfliteral{0 Tr }%
  \setbox\z@=\hbox{%
    \begin{tikzpicture}[baseline=(gtnode.base)]%
      \node[inner sep=0pt](gtnode){\phantom{#3}};%
      \shade[left color=#1, right color=#2]%
        (gtnode.south west)rectangle(gtnode.north east);%
    \end{tikzpicture}%
  }%
  \dimen@=\dimexpr\wd\z@ * 7200 / 7227\relax%
  \edef\gt@picwd{\strip@pt\dimen@}%
  \box\z@%
  \pdfliteral{Q}%
  \pdfliteral{1 0 0 1 \gt@picwd\space 0 cm}%
}
\makeatother

% ── Dégradés nommés (utilisés par \diagbg et \layeredbg) ─────────────────────
% sweep4 : froid → chaud (bleu / turquoise / orange / magenta)
\pgfdeclarehorizontalshading{sweep4}{100bp}{%
  color(0bp)=(C1lo);
  color(33bp)=(C1hi);
  color(66bp)=(C2lo);
  color(100bp)=(C2hi)%
}
% sweep3 : bleu → vert → orange
\pgfdeclarehorizontalshading{sweep3}{100bp}{%
  color(0bp)=(C1lo);
  color(50bp)=(C3lo);
  color(100bp)=(C2lo)%
}
% nightdawn : sombre → bleu → turquoise → orange
\pgfdeclarehorizontalshading{nightdawn}{100bp}{%
  color(0bp)=(Cdark);
  color(30bp)=(C1lo);
  color(60bp)=(C1hi);
  color(100bp)=(C2lo)%
}
% dusk : magenta → bleu → turquoise
\pgfdeclarehorizontalshading{dusk}{100bp}{%
  color(0bp)=(C2hi);
  color(45bp)=(C1lo);
  color(100bp)=(C1hi)%
}

% ── Anneaux radiaux (outer = white : fond blanc obligatoire) ──────────────────
\pgfdeclareradialshading{ringsA}{\pgfpointorigin}{%
  color(0.0cm)=(white!20!C1hi);
  color(0.9cm)=(C1hi);
  color(2.0cm)=(C1lo);
  color(3.4cm)=(C1lo!30!gray);
  color(4.8cm)=(gray!45);
  color(6.2cm)=(gray!18);
  color(7.6cm)=(gray!6);
  color(9.0cm)=(white)%
}
\pgfdeclareradialshading{ringsB}{\pgfpointorigin}{%
  color(0.0cm)=(white!20!C2hi);
  color(0.9cm)=(C2hi);
  color(2.0cm)=(C2lo);
  color(3.4cm)=(C2lo!30!gray);
  color(4.8cm)=(gray!45);
  color(6.2cm)=(gray!18);
  color(7.6cm)=(gray!6);
  color(9.0cm)=(white)%
}
\pgfdeclareradialshading{ringsC}{\pgfpointorigin}{%
  color(0.0cm)=(white!20!C3hi);
  color(0.9cm)=(C3hi);
  color(2.0cm)=(C3lo);
  color(3.4cm)=(C3lo!30!gray);
  color(4.8cm)=(gray!45);
  color(6.2cm)=(gray!18);
  color(7.6cm)=(gray!6);
  color(9.0cm)=(white)%
}

% ── Fonds de page ─────────────────────────────────────────────────────────────

% \placerings{ringsA|B|C}{anchor}{xcm}{ycm}
\newcommand{\placerings}[4]{%
  \AddToShipoutPictureBG*{%
    \begin{tikzpicture}[remember picture, overlay]
      \begin{scope}
        \clip (current page.south west) rectangle (current page.north east);
        \begin{scope}[shift={([xshift=#3 cm, yshift=#4 cm]current page.#2)}]
          \pgfuseshading{#1}
        \end{scope}
      \end{scope}
    \end{tikzpicture}%
  }%
}

% \spotfade{color}{xcm}{ycm}{anchor}{maxradius}
% À appeler à l'intérieur d'un tikzpicture remember picture/overlay.
\newcommand{\spotfade}[5]{%
  \begin{scope}
    \clip (current page.south west) rectangle (current page.north east);
    \pgfmathsetmacro{\Rmax}{#5}
    \foreach \frac/\op in {%
        1.00/0.04, 0.88/0.09, 0.76/0.14,
        0.64/0.20, 0.52/0.28, 0.40/0.38,
        0.28/0.49, 0.18/0.61, 0.10/0.72, 0.04/0.82}{%
      \pgfmathsetmacro{\r}{\Rmax*\frac}
      \fill[#1, opacity=\op]
        ([xshift=#2 cm, yshift=#3 cm]current page.#4) circle(\r cm);
    }
  \end{scope}
}

% \complexbg{baselo}{basehi}{spotcolor}{xcm}{ycm}{anchor}{radius}
\newcommand{\complexbg}[7]{%
  \AddToShipoutPictureBG*{%
    \begin{tikzpicture}[remember picture, overlay]
      \shade[top color=#1, bottom color=#2]
        (current page.south west) rectangle (current page.north east);
      \spotfade{#3}{#4}{#5}{#6}{#7}
    \end{tikzpicture}%
  }%
}

% \fourcornerbg{SW}{NW}{SE}{NE}
\newcommand{\fourcornerbg}[4]{%
  \AddToShipoutPictureBG*{%
    \begin{tikzpicture}[remember picture, overlay]
      \shade[lower left=#1, upper left=#2, lower right=#3, upper right=#4]
        (current page.south west) rectangle (current page.north east);
    \end{tikzpicture}%
  }%
}

% \halftopbg{SW}{NW}{SE}{NE}  — moitié haute colorée (bilinéaire), bas blanc
% Liseré solide en bas de page reprenant le dégradé gauche/droite du fond.
\newcommand{\halftopbg}[4]{%
  \AddToShipoutPictureBG*{%
    \begin{tikzpicture}[remember picture, overlay]
      \shade[lower left=#1, upper left=#2, lower right=#3, upper right=#4]
        (current page.west) rectangle (current page.north east);
      \shade[left color=#1, right color=#3]
        (current page.south west)
        rectangle
        ([yshift=12pt]current page.south east);
    \end{tikzpicture}%
  }%
}

% \topbg{splitht}{SW}{NW}{SE}{NE}  — haut coloré jusqu'à splitht depuis bas papier
% splitht : hauteur depuis current page.south (ex. 129.4mm pour alignement custom)
\newcommand{\topbg}[5]{%
  \AddToShipoutPictureBG*{%
    \begin{tikzpicture}[remember picture, overlay]
      \shade[lower left=#2, upper left=#3, lower right=#4, upper right=#5]
        ([yshift=#1]current page.south west)
        rectangle
        (current page.north east);
      \shade[left color=#2, right color=#4]
        (current page.south west)
        rectangle
        ([yshift=12pt]current page.south east);
    \end{tikzpicture}%
  }%
}

% \topbgsplit{splitht}{SW}{NW}{SE}{NE}  — comme \topbg, sans liseré bas (dessous blanc)
% splitht peut être \zposy{label}sp pour un split dynamique via zref-savepos.
\newcommand{\topbgsplit}[5]{%
  \AddToShipoutPictureBG*{%
    \begin{tikzpicture}[remember picture, overlay]
      \shade[lower left=#2, upper left=#3, lower right=#4, upper right=#5]
        ([yshift=#1]current page.south west)
        rectangle
        (current page.north east);
    \end{tikzpicture}%
  }%
}

% \halfbotbg{SW}{NW}{SE}{NE}  — moitié basse colorée, haut blanc
\newcommand{\halfbotbg}[4]{%
  \AddToShipoutPictureBG*{%
    \begin{tikzpicture}[remember picture, overlay]
      \shade[lower left=#1, upper left=#2, lower right=#3, upper right=#4]
        (current page.south west) rectangle (current page.east);
    \end{tikzpicture}%
  }%
}

% \diagbg{sweep4|sweep3|nightdawn|dusk}{angle}
%   angle : 0 (→)  45 (↗)  90 (↑)  135 (↖)  180 (←)
\newcommand{\diagbg}[2]{%
  \AddToShipoutPictureBG*{%
    \begin{tikzpicture}[remember picture, overlay]
      \shade[shading=#1, shading angle=#2]
        (current page.south west) rectangle (current page.north east);
    \end{tikzpicture}%
  }%
}

% \layeredbg{shading1}{angle1}{shading2}{angle2}{opacity2}
% Couche 2 à opacity2 par-dessus couche 1 : angles croisés → 4 coins distincts.
\newcommand{\layeredbg}[5]{%
  \AddToShipoutPictureBG*{%
    \begin{tikzpicture}[remember picture, overlay]
      \shade[shading=#1, shading angle=#2]
        (current page.south west) rectangle (current page.north east);
      \begin{scope}[opacity=#5]
        \shade[shading=#3, shading angle=#4]
          (current page.south west) rectangle (current page.north east);
      \end{scope}
    \end{tikzpicture}%
  }%
}

% \fourcornerspot{SW}{NW}{SE}{NE}{spotcolor}{radius}
\newcommand{\fourcornerspot}[6]{%
  \AddToShipoutPictureBG*{%
    \begin{tikzpicture}[remember picture, overlay]
      \shade[lower left=#1, upper left=#2, lower right=#3, upper right=#4]
        (current page.south west) rectangle (current page.north east);
      \spotfade{#5}{0}{0}{center}{#6}
    \end{tikzpicture}%
  }%
}

% \trispot{c-NW}{c-NE}{c-S}   fond sombre + halos NW / NE / bas-centre
\newcommand{\trispot}[3]{%
  \AddToShipoutPictureBG*{%
    \begin{tikzpicture}[remember picture, overlay]
      \fill[Cdark] (current page.south west) rectangle (current page.north east);
      \spotfade{#1}{0}{0}{north west}{12}
      \spotfade{#2}{0}{0}{north east}{12}
      \spotfade{#3}{0}{0}{south}{10}
    \end{tikzpicture}%
  }%
}

% \tricorner{c-NW}{c-NE}{c-SE}   fond sombre + halos NW / NE / SE
\newcommand{\tricorner}[3]{%
  \AddToShipoutPictureBG*{%
    \begin{tikzpicture}[remember picture, overlay]
      \fill[Cdark] (current page.south west) rectangle (current page.north east);
      \spotfade{#1}{0}{0}{north west}{12}
      \spotfade{#2}{0}{0}{north east}{12}
      \spotfade{#3}{0}{0}{south east}{12}
    \end{tikzpicture}%
  }%
}

% \tricornergb{c-NW}{c-NE}{c-SE}   gradient bilinéaire 3 coins, SW fixé à Cdark
\newcommand{\tricornergb}[3]{\fourcornerbg{Cdark}{#1}{#3}{#2}}

% \trispotgb{c-NW}{c-NE}{c-bas}   gradient bilinéaire, coins bas identiques
\newcommand{\trispotgb}[3]{\fourcornerbg{#3}{#1}{#3}{#2}}

% LaTeX limite à 9 args/commande : \doublebgA stocke les params, \doublebgB déclenche le rendu.
% \doublebgA{baselo}{basehi}{spot1color}{s1x}{s1y}{s1anchor}{s1radius}
% \doublebgB{spot2color}{s2x}{s2y}{s2anchor}{s2radius}
\newcommand{\doublebgA}[7]{%
  \gdef\dbblo{#1}\gdef\dbbhi{#2}%
  \gdef\dbca{#3}\gdef\dbxa{#4}\gdef\dbya{#5}\gdef\dbancha{#6}\gdef\dbra{#7}%
}
\newcommand{\doublebgB}[5]{%
  \AddToShipoutPictureBG*{%
    \begin{tikzpicture}[remember picture, overlay]
      \shade[top color=\dbblo, bottom color=\dbbhi]
        (current page.south west) rectangle (current page.north east);
      \spotfade{\dbca}{\dbxa}{\dbya}{\dbancha}{\dbra}
      \spotfade{#1}{#2}{#3}{#4}{#5}
    \end{tikzpicture}%
  }%
}

% ── Boites tcolorbox ─────────────────────────────────────────────────────────
\usepackage[most]{tcolorbox}
\tcbuselibrary{skins, breakable}

\newtcolorbox{gradbox}[2]{
  enhanced, breakable, arc=0pt, boxrule=0pt,
  interior style={left color=#1, right color=#2}, coltext=white,
  left=1.5cm, right=1.5cm, top=1cm, bottom=1cm,
  before skip=0.9em, after skip=0.9em,
}
\newtcolorbox{fullband}[2]{
  enhanced, breakable, arc=0pt, boxrule=0pt,
  interior style={left color=#1, right color=#2}, coltext=white,
  enlarge left by=-2.5cm, enlarge right by=-2.5cm,
  left=4.1cm, right=4.1cm, top=1cm, bottom=1cm,
  before skip=0.9em, after skip=0.9em,
}
\newtcolorbox{darkbox}{
  enhanced, breakable, arc=0pt, boxrule=0pt,
  colback=Cdark, colframe=Cdark, coltext=white,
  left=1.5cm, right=1.5cm, top=1cm, bottom=1cm,
  before skip=0.9em, after skip=0.9em,
}
\newtcolorbox{lightbox}{
  enhanced, breakable, arc=3pt, boxrule=0pt,
  colback=Clight, colframe=Clight,
  left=1.2cm, right=1.2cm, top=0.8cm, bottom=0.8cm,
  before skip=0.9em, after skip=0.9em,
}
\newtcolorbox{solidbox}[1]{
  enhanced, breakable, arc=0pt, boxrule=0pt,
  colback=#1, colframe=#1, coltext=white,
  left=1cm, right=1cm, top=0.8cm, bottom=0.8cm,
}
\newtcolorbox{whitebox}{
  enhanced, breakable, arc=3pt, boxrule=0pt,
  colback=white, colframe=white, coltext=black,
  drop shadow,
  left=1.2cm, right=1.2cm, top=0.9cm, bottom=0.9cm,
  before skip=0.9em, after skip=0.9em,
}
\newtcolorbox{codebox}{
  enhanced, breakable, arc=3pt, boxrule=0.4pt,
  colback=Ccode, colframe=gray!35, coltext=black,
  left=0.9cm, right=0.9cm, top=0.6cm, bottom=0.6cm,
  fontupper=\small\ttfamily,
  before skip=0.9em, after skip=0.9em,
}

\newcommand{\nbullet}[2]{%
  \begin{tikzpicture}[baseline=(c.base)]
    \node[circle, minimum size=1.6em, inner sep=0pt,
          top color=C1lo, bottom color=C1hi,
          text=white, font=\bfseries\small] (c) {#1};
  \end{tikzpicture}~~#2%
}

% ── Titres ───────────────────────────────────────────────────────────────────
\usepackage{titlesec}
\titleformat{\section}  {\titlefont\bfseries\Huge\raggedright}  {}{0pt}{}
\titleformat{\subsection}{\titlefont\bfseries\Large\raggedright}{}{0pt}{}
\titlespacing{\section}  {0pt}{1.5em}{1em}
\titlespacing{\subsection}{0pt}{1em}{0.3em}

\newcommand{\seclabel}[2]{%
  {\small\titlefont\bfseries\color{#1}\MakeUppercase{#2}}\par\vspace{0.8em}%
}
\newcommand{\param}[1]{\textcolor{C2lo}{\texttt{#1}}}

% ── Fonds de page (eso-pic) ───────────────────────────────────────────────────
\usepackage{eso-pic}
\usepackage{pdflscape}

% Requiert : \documentclass déjà déclaré

% ── Police ───────────────────────────────────────────────────────────────────
\usepackage[T1]{fontenc}
\usepackage[utf8]{inputenc}
\usepackage{tgheros}
\renewcommand{\familydefault}{\sfdefault}
\newcommand{\titlefont}{\sffamily}

% ── Mise en page ──────────────────────────────────────────────────────────────
\usepackage[a4paper, top=2cm, bottom=2.2cm, left=2.5cm, right=2.5cm]{geometry}
\usepackage{microtype}
\setlength{\parindent}{0pt}
\setlength{\parskip}{0.6em}
\pagestyle{empty}

% ── Couleurs ─────────────────────────────────────────────────────────────────
\usepackage{xcolor}
\definecolor{C1lo}{HTML}{1E3A8A}  % bleu foncé
\definecolor{C1hi}{HTML}{0D9488}  % turquoise
\definecolor{C2lo}{HTML}{F4801A}  % orange
\definecolor{C2hi}{HTML}{C2185B}  % magenta
\definecolor{C3lo}{HTML}{27AE60}  % vert
\definecolor{C3hi}{HTML}{0D9488}  % turquoise (alias C1hi)
\definecolor{Cdark}{HTML}{1A1A2E}
\definecolor{Clight}{HTML}{F3F4F6}
\definecolor{Ccode}{HTML}{F0F0F0}

\definecolor{Cgrey}{HTML}{868686}
\colorlet{grey}{Cgrey}
\definecolor{Cc1}{HTML}{F0F0F0}
\definecolor{Cc2}{HTML}{F0F0F0}

% ── Mise en page avancée ─────────────────────────────────────────────────────
\usepackage{tabularx}
\usepackage{calc}
\usepackage{wrapfig}

% ── Commandes stat ────────────────────────────────────────────────────────────
%
% Six commandes de mise en page côte-à-côte : grand chiffre + sous-titre
% dans une colonne (gauche ou droite), texte courant dans l'autre colonne.
%
% Principe commun (variantes vtop/hbox)
% ──────────────────────────────────────
%   \hbox to \linewidth{ vtop(stat) vtop(texte) \hss }
%
%   vtop avec \kern0pt : point de référence au sommet (height=0, depth=contenu).
%   \parskip=0pt dans le vtop stat empêche un saut parasite avant le gros chiffre.
%
%   Espacement — variable \statgap
%   ──────────────────────────────
%   \statgap = blanc visible entre le chiffre et le texte courant.
%   Commun aux 6 commandes : changer \statgap recalibre tout.
%   Implémentation vtop/hbox : le blanc est absorbé dans la largeur du vtop-stat
%   (statwidth = natural + \statgap, pas de \hspace explicite entre les vtops).
%   Implémentation wrapfigure : \columnsep = \statgap (gap géré par wrapfig).
%
%   Alignement optique (sommet des caractères)
%   ───────────────────────────────────────────
%   Le vtop-texte est relevé par \raise\dimexpr\statcap-\bodycap\relax pour
%   aligner les sommets de glyphes. Non appliqué aux variantes wrapfigure.
%
% Comportement d'ancrage
% ──────────────────────
%   Les variantes vtop/hbox produisent un \hbox nu. Le caller gère \par et blancs.
%
% Les six commandes
% ─────────────────
%   \statcol    stat GAUCHE, vtop/hbox
%   \statflow   stat GAUCHE, vtop/hbox  (alias fonctionnel de \statcol)
%   \statflowr  stat DROITE, vtop/hbox, raggedleft, chiffre à la marge droite
%   \statcolwf  stat GAUCHE, wrapfigure (float : hors vbox/minipage)
%   \statflowwf stat GAUCHE, wrapfigure (idem, alias de \statcolwf)
%   \statflowrwf stat DROITE, wrapfigure (idem)

\newlength{\statwidth}
\newlength{\statcap}
\newlength{\bodycap}
\newlength{\statgap}
\setlength{\statgap}{1.0cm}

% \statcol{chiffre}{sous-titre}{texte}
\newcommand{\statcol}[3]{%
  \settowidth{\statwidth}{{\fontsize{60}{64}\selectfont\bfseries #1}}%
  \addtolength{\statwidth}{\statgap}%
  \settoheight{\statcap}{{\fontsize{60}{64}\selectfont\bfseries X}}%
  \settoheight{\bodycap}{{\large X}}%
  \noindent\hbox to \linewidth{%
    \vtop{\kern0pt\hsize=\statwidth\relax
      \parskip=0pt\relax
      {\fontsize{60}{64}\selectfont\bfseries #1\par}%
      {\large #2}%
    }%
    \raise\dimexpr\statcap-\bodycap\relax
    \vtop{\kern0pt\hsize=\dimexpr\linewidth-\statwidth\relax
      {\large #3\par}%
    }%
    \hss
  }%
}

% \statflow{chiffre}{sous-titre}{texte}  — stat gauche
\newcommand{\statflow}[3]{%
  \settowidth{\statwidth}{{\fontsize{60}{64}\selectfont\bfseries #1}}%
  \addtolength{\statwidth}{\statgap}%
  \settoheight{\statcap}{{\fontsize{60}{64}\selectfont\bfseries X}}%
  \settoheight{\bodycap}{{\large X}}%
  \noindent\hbox to \linewidth{%
    \vtop{\kern0pt\hsize=\statwidth\relax
      \parskip=0pt\relax
      {\fontsize{60}{64}\selectfont\bfseries #1\par}%
      {\large #2}%
    }%
    \raise\dimexpr\statcap-\bodycap\relax
    \vtop{\kern0pt\hsize=\dimexpr\linewidth-\statwidth\relax
      {\large #3\par}%
    }%
    \hss
  }%
}

% \statflowr{chiffre}{sous-titre}{texte}  — stat droite
\newcommand{\statflowr}[3]{%
  \settowidth{\statwidth}{{\fontsize{60}{64}\selectfont\bfseries #1}}%
  \addtolength{\statwidth}{\statgap}%
  \settoheight{\statcap}{{\fontsize{60}{64}\selectfont\bfseries X}}%
  \settoheight{\bodycap}{{\large X}}%
  \noindent\hbox to \linewidth{%
    \raise\dimexpr\statcap-\bodycap\relax
    \vtop{\kern0pt\hsize=\dimexpr\linewidth-\statwidth\relax
      {\large #3\par}%
    }%
    \vtop{\kern0pt\hsize=\statwidth\relax
      \parskip=0pt\relax
      \raggedleft
      {\fontsize{60}{64}\selectfont\bfseries #1\par}%
      {\large #2\par}%
    }%
    \hss
  }%
}

% ── Variantes wrapfigure ──────────────────────────────────────────────────────
%
% Doublons de statcol/statflow/statflowr utilisant wrapfigure.
% Le texte (#3) s'écoule autour du chiffre ; le contenu qui suit la commande
% dans le document wrappera aussi (comportement natif wrapfigure).
%
% Contrainte : wrapfigure est un float. Il ne fonctionne pas dans un vbox,
% minipage ou toute boîte : le float est alors différé hors de la boîte.
% À réserver au flot normal du document.
%
% Gap : \columnsep mis à \statgap en local.
% \intextsep mis à 0pt pour ne pas creuser de vide au-dessus du chiffre.
% La largeur déclarée à wrapfigure = natural_width du chiffre seul (le gap
% est pris en charge par \columnsep, pas inclus dans la boîte figure).

% \statcolwf{chiffre}{sous-titre}{texte}  — stat gauche, wrapfigure
\newcommand{\statcolwf}[3]{%
  \settowidth{\statwidth}{{\fontsize{60}{64}\selectfont\bfseries #1}}%
  \setlength{\intextsep}{0pt}%
  \setlength{\columnsep}{\statgap}%
  \begin{wrapfigure}{l}{\statwidth}%
    \vtop{\kern0pt\hsize=\statwidth\relax
      \parskip=0pt\relax
      {\fontsize{60}{64}\selectfont\bfseries #1\par}%
      {\large #2}%
    }%
  \end{wrapfigure}%
  {\large #3\par}%
}

% \statflowwf{chiffre}{sous-titre}{texte}  — stat gauche, wrapfigure
\newcommand{\statflowwf}[3]{%
  \settowidth{\statwidth}{{\fontsize{60}{64}\selectfont\bfseries #1}}%
  \setlength{\intextsep}{0pt}%
  \setlength{\columnsep}{\statgap}%
  \begin{wrapfigure}{l}{\statwidth}%
    \vtop{\kern0pt\hsize=\statwidth\relax
      \parskip=0pt\relax
      {\fontsize{60}{64}\selectfont\bfseries #1\par}%
      {\large #2}%
    }%
  \end{wrapfigure}%
  {\large #3\par}%
}

% \statflowrwf{chiffre}{sous-titre}{texte}  — stat droite, wrapfigure
\newcommand{\statflowrwf}[3]{%
  \settowidth{\statwidth}{{\fontsize{60}{64}\selectfont\bfseries #1}}%
  \setlength{\intextsep}{0pt}%
  \setlength{\columnsep}{\statgap}%
  \begin{wrapfigure}{r}{\statwidth}%
    \vtop{\kern0pt\hsize=\statwidth\relax
      \parskip=0pt\relax
      \raggedleft
      {\fontsize{60}{64}\selectfont\bfseries #1\par}%
      {\large #2\par}%
    }%
  \end{wrapfigure}%
  {\large #3\par}%
}

% ── TikZ ─────────────────────────────────────────────────────────────────────
\usepackage{tikz}
\usetikzlibrary{shadings, shadows, fadings}

% \gradtext{couleur-gauche}{couleur-droite}{contenu}
% Texte en dégradé gauche → droite, sur n'importe quel fond.
% Méthode : PDF Tr=7 (clip aux contours des glyphes) + dégradé TikZ.
% Après Q, le tracker CTM de pdflatex est désynchronisé ; on réémet un cm
% compensatoire égal à la largeur de l'image TikZ (convertie en bp).
\makeatletter
\newcommand{\gradtext}[3]{%
  \leavevmode%
  \pdfliteral{q 7 Tr }%
  \rlap{{#3}}%
  \pdfliteral{0 Tr }%
  \setbox\z@=\hbox{%
    \begin{tikzpicture}[baseline=(gtnode.base)]%
      \node[inner sep=0pt](gtnode){\phantom{#3}};%
      \shade[left color=#1, right color=#2]%
        (gtnode.south west)rectangle(gtnode.north east);%
    \end{tikzpicture}%
  }%
  \dimen@=\dimexpr\wd\z@ * 7200 / 7227\relax%
  \edef\gt@picwd{\strip@pt\dimen@}%
  \box\z@%
  \pdfliteral{Q}%
  \pdfliteral{1 0 0 1 \gt@picwd\space 0 cm}%
}
\makeatother

% ── Dégradés nommés (utilisés par \diagbg et \layeredbg) ─────────────────────
% sweep4 : froid → chaud (bleu / turquoise / orange / magenta)
\pgfdeclarehorizontalshading{sweep4}{100bp}{%
  color(0bp)=(C1lo);
  color(33bp)=(C1hi);
  color(66bp)=(C2lo);
  color(100bp)=(C2hi)%
}
% sweep3 : bleu → vert → orange
\pgfdeclarehorizontalshading{sweep3}{100bp}{%
  color(0bp)=(C1lo);
  color(50bp)=(C3lo);
  color(100bp)=(C2lo)%
}
% nightdawn : sombre → bleu → turquoise → orange
\pgfdeclarehorizontalshading{nightdawn}{100bp}{%
  color(0bp)=(Cdark);
  color(30bp)=(C1lo);
  color(60bp)=(C1hi);
  color(100bp)=(C2lo)%
}
% dusk : magenta → bleu → turquoise
\pgfdeclarehorizontalshading{dusk}{100bp}{%
  color(0bp)=(C2hi);
  color(45bp)=(C1lo);
  color(100bp)=(C1hi)%
}

% ── Anneaux radiaux (outer = white : fond blanc obligatoire) ──────────────────
\pgfdeclareradialshading{ringsA}{\pgfpointorigin}{%
  color(0.0cm)=(white!20!C1hi);
  color(0.9cm)=(C1hi);
  color(2.0cm)=(C1lo);
  color(3.4cm)=(C1lo!30!gray);
  color(4.8cm)=(gray!45);
  color(6.2cm)=(gray!18);
  color(7.6cm)=(gray!6);
  color(9.0cm)=(white)%
}
\pgfdeclareradialshading{ringsB}{\pgfpointorigin}{%
  color(0.0cm)=(white!20!C2hi);
  color(0.9cm)=(C2hi);
  color(2.0cm)=(C2lo);
  color(3.4cm)=(C2lo!30!gray);
  color(4.8cm)=(gray!45);
  color(6.2cm)=(gray!18);
  color(7.6cm)=(gray!6);
  color(9.0cm)=(white)%
}
\pgfdeclareradialshading{ringsC}{\pgfpointorigin}{%
  color(0.0cm)=(white!20!C3hi);
  color(0.9cm)=(C3hi);
  color(2.0cm)=(C3lo);
  color(3.4cm)=(C3lo!30!gray);
  color(4.8cm)=(gray!45);
  color(6.2cm)=(gray!18);
  color(7.6cm)=(gray!6);
  color(9.0cm)=(white)%
}

% ── Fonds de page ─────────────────────────────────────────────────────────────

% \placerings{ringsA|B|C}{anchor}{xcm}{ycm}
\newcommand{\placerings}[4]{%
  \AddToShipoutPictureBG*{%
    \begin{tikzpicture}[remember picture, overlay]
      \begin{scope}
        \clip (current page.south west) rectangle (current page.north east);
        \begin{scope}[shift={([xshift=#3 cm, yshift=#4 cm]current page.#2)}]
          \pgfuseshading{#1}
        \end{scope}
      \end{scope}
    \end{tikzpicture}%
  }%
}

% \spotfade{color}{xcm}{ycm}{anchor}{maxradius}
% À appeler à l'intérieur d'un tikzpicture remember picture/overlay.
\newcommand{\spotfade}[5]{%
  \begin{scope}
    \clip (current page.south west) rectangle (current page.north east);
    \pgfmathsetmacro{\Rmax}{#5}
    \foreach \frac/\op in {%
        1.00/0.04, 0.88/0.09, 0.76/0.14,
        0.64/0.20, 0.52/0.28, 0.40/0.38,
        0.28/0.49, 0.18/0.61, 0.10/0.72, 0.04/0.82}{%
      \pgfmathsetmacro{\r}{\Rmax*\frac}
      \fill[#1, opacity=\op]
        ([xshift=#2 cm, yshift=#3 cm]current page.#4) circle(\r cm);
    }
  \end{scope}
}

% \complexbg{baselo}{basehi}{spotcolor}{xcm}{ycm}{anchor}{radius}
\newcommand{\complexbg}[7]{%
  \AddToShipoutPictureBG*{%
    \begin{tikzpicture}[remember picture, overlay]
      \shade[top color=#1, bottom color=#2]
        (current page.south west) rectangle (current page.north east);
      \spotfade{#3}{#4}{#5}{#6}{#7}
    \end{tikzpicture}%
  }%
}

% \fourcornerbg{SW}{NW}{SE}{NE}
\newcommand{\fourcornerbg}[4]{%
  \AddToShipoutPictureBG*{%
    \begin{tikzpicture}[remember picture, overlay]
      \shade[lower left=#1, upper left=#2, lower right=#3, upper right=#4]
        (current page.south west) rectangle (current page.north east);
    \end{tikzpicture}%
  }%
}

% \halftopbg{SW}{NW}{SE}{NE}  — moitié haute colorée (bilinéaire), bas blanc
% Liseré solide en bas de page reprenant le dégradé gauche/droite du fond.
\newcommand{\halftopbg}[4]{%
  \AddToShipoutPictureBG*{%
    \begin{tikzpicture}[remember picture, overlay]
      \shade[lower left=#1, upper left=#2, lower right=#3, upper right=#4]
        (current page.west) rectangle (current page.north east);
      \shade[left color=#1, right color=#3]
        (current page.south west)
        rectangle
        ([yshift=12pt]current page.south east);
    \end{tikzpicture}%
  }%
}

% \topbg{splitht}{SW}{NW}{SE}{NE}  — haut coloré jusqu'à splitht depuis bas papier
% splitht : hauteur depuis current page.south (ex. 129.4mm pour alignement custom)
\newcommand{\topbg}[5]{%
  \AddToShipoutPictureBG*{%
    \begin{tikzpicture}[remember picture, overlay]
      \shade[lower left=#2, upper left=#3, lower right=#4, upper right=#5]
        ([yshift=#1]current page.south west)
        rectangle
        (current page.north east);
      \shade[left color=#2, right color=#4]
        (current page.south west)
        rectangle
        ([yshift=12pt]current page.south east);
    \end{tikzpicture}%
  }%
}

% \topbgsplit{splitht}{SW}{NW}{SE}{NE}  — comme \topbg, sans liseré bas (dessous blanc)
% splitht peut être \zposy{label}sp pour un split dynamique via zref-savepos.
\newcommand{\topbgsplit}[5]{%
  \AddToShipoutPictureBG*{%
    \begin{tikzpicture}[remember picture, overlay]
      \shade[lower left=#2, upper left=#3, lower right=#4, upper right=#5]
        ([yshift=#1]current page.south west)
        rectangle
        (current page.north east);
    \end{tikzpicture}%
  }%
}

% \halfbotbg{SW}{NW}{SE}{NE}  — moitié basse colorée, haut blanc
\newcommand{\halfbotbg}[4]{%
  \AddToShipoutPictureBG*{%
    \begin{tikzpicture}[remember picture, overlay]
      \shade[lower left=#1, upper left=#2, lower right=#3, upper right=#4]
        (current page.south west) rectangle (current page.east);
    \end{tikzpicture}%
  }%
}

% \diagbg{sweep4|sweep3|nightdawn|dusk}{angle}
%   angle : 0 (→)  45 (↗)  90 (↑)  135 (↖)  180 (←)
\newcommand{\diagbg}[2]{%
  \AddToShipoutPictureBG*{%
    \begin{tikzpicture}[remember picture, overlay]
      \shade[shading=#1, shading angle=#2]
        (current page.south west) rectangle (current page.north east);
    \end{tikzpicture}%
  }%
}

% \layeredbg{shading1}{angle1}{shading2}{angle2}{opacity2}
% Couche 2 à opacity2 par-dessus couche 1 : angles croisés → 4 coins distincts.
\newcommand{\layeredbg}[5]{%
  \AddToShipoutPictureBG*{%
    \begin{tikzpicture}[remember picture, overlay]
      \shade[shading=#1, shading angle=#2]
        (current page.south west) rectangle (current page.north east);
      \begin{scope}[opacity=#5]
        \shade[shading=#3, shading angle=#4]
          (current page.south west) rectangle (current page.north east);
      \end{scope}
    \end{tikzpicture}%
  }%
}

% \fourcornerspot{SW}{NW}{SE}{NE}{spotcolor}{radius}
\newcommand{\fourcornerspot}[6]{%
  \AddToShipoutPictureBG*{%
    \begin{tikzpicture}[remember picture, overlay]
      \shade[lower left=#1, upper left=#2, lower right=#3, upper right=#4]
        (current page.south west) rectangle (current page.north east);
      \spotfade{#5}{0}{0}{center}{#6}
    \end{tikzpicture}%
  }%
}

% \trispot{c-NW}{c-NE}{c-S}   fond sombre + halos NW / NE / bas-centre
\newcommand{\trispot}[3]{%
  \AddToShipoutPictureBG*{%
    \begin{tikzpicture}[remember picture, overlay]
      \fill[Cdark] (current page.south west) rectangle (current page.north east);
      \spotfade{#1}{0}{0}{north west}{12}
      \spotfade{#2}{0}{0}{north east}{12}
      \spotfade{#3}{0}{0}{south}{10}
    \end{tikzpicture}%
  }%
}

% \tricorner{c-NW}{c-NE}{c-SE}   fond sombre + halos NW / NE / SE
\newcommand{\tricorner}[3]{%
  \AddToShipoutPictureBG*{%
    \begin{tikzpicture}[remember picture, overlay]
      \fill[Cdark] (current page.south west) rectangle (current page.north east);
      \spotfade{#1}{0}{0}{north west}{12}
      \spotfade{#2}{0}{0}{north east}{12}
      \spotfade{#3}{0}{0}{south east}{12}
    \end{tikzpicture}%
  }%
}

% \tricornergb{c-NW}{c-NE}{c-SE}   gradient bilinéaire 3 coins, SW fixé à Cdark
\newcommand{\tricornergb}[3]{\fourcornerbg{Cdark}{#1}{#3}{#2}}

% \trispotgb{c-NW}{c-NE}{c-bas}   gradient bilinéaire, coins bas identiques
\newcommand{\trispotgb}[3]{\fourcornerbg{#3}{#1}{#3}{#2}}

% LaTeX limite à 9 args/commande : \doublebgA stocke les params, \doublebgB déclenche le rendu.
% \doublebgA{baselo}{basehi}{spot1color}{s1x}{s1y}{s1anchor}{s1radius}
% \doublebgB{spot2color}{s2x}{s2y}{s2anchor}{s2radius}
\newcommand{\doublebgA}[7]{%
  \gdef\dbblo{#1}\gdef\dbbhi{#2}%
  \gdef\dbca{#3}\gdef\dbxa{#4}\gdef\dbya{#5}\gdef\dbancha{#6}\gdef\dbra{#7}%
}
\newcommand{\doublebgB}[5]{%
  \AddToShipoutPictureBG*{%
    \begin{tikzpicture}[remember picture, overlay]
      \shade[top color=\dbblo, bottom color=\dbbhi]
        (current page.south west) rectangle (current page.north east);
      \spotfade{\dbca}{\dbxa}{\dbya}{\dbancha}{\dbra}
      \spotfade{#1}{#2}{#3}{#4}{#5}
    \end{tikzpicture}%
  }%
}

% ── Boites tcolorbox ─────────────────────────────────────────────────────────
\usepackage[most]{tcolorbox}
\tcbuselibrary{skins, breakable}

\newtcolorbox{gradbox}[2]{
  enhanced, breakable, arc=0pt, boxrule=0pt,
  interior style={left color=#1, right color=#2}, coltext=white,
  left=1.5cm, right=1.5cm, top=1cm, bottom=1cm,
  before skip=0.9em, after skip=0.9em,
}
\newtcolorbox{fullband}[2]{
  enhanced, breakable, arc=0pt, boxrule=0pt,
  interior style={left color=#1, right color=#2}, coltext=white,
  enlarge left by=-2.5cm, enlarge right by=-2.5cm,
  left=4.1cm, right=4.1cm, top=1cm, bottom=1cm,
  before skip=0.9em, after skip=0.9em,
}
\newtcolorbox{darkbox}{
  enhanced, breakable, arc=0pt, boxrule=0pt,
  colback=Cdark, colframe=Cdark, coltext=white,
  left=1.5cm, right=1.5cm, top=1cm, bottom=1cm,
  before skip=0.9em, after skip=0.9em,
}
\newtcolorbox{lightbox}{
  enhanced, breakable, arc=3pt, boxrule=0pt,
  colback=Clight, colframe=Clight,
  left=1.2cm, right=1.2cm, top=0.8cm, bottom=0.8cm,
  before skip=0.9em, after skip=0.9em,
}
\newtcolorbox{solidbox}[1]{
  enhanced, breakable, arc=0pt, boxrule=0pt,
  colback=#1, colframe=#1, coltext=white,
  left=1cm, right=1cm, top=0.8cm, bottom=0.8cm,
}
\newtcolorbox{whitebox}{
  enhanced, breakable, arc=3pt, boxrule=0pt,
  colback=white, colframe=white, coltext=black,
  drop shadow,
  left=1.2cm, right=1.2cm, top=0.9cm, bottom=0.9cm,
  before skip=0.9em, after skip=0.9em,
}
\newtcolorbox{codebox}{
  enhanced, breakable, arc=3pt, boxrule=0.4pt,
  colback=Ccode, colframe=gray!35, coltext=black,
  left=0.9cm, right=0.9cm, top=0.6cm, bottom=0.6cm,
  fontupper=\small\ttfamily,
  before skip=0.9em, after skip=0.9em,
}

\newcommand{\nbullet}[2]{%
  \begin{tikzpicture}[baseline=(c.base)]
    \node[circle, minimum size=1.6em, inner sep=0pt,
          top color=C1lo, bottom color=C1hi,
          text=white, font=\bfseries\small] (c) {#1};
  \end{tikzpicture}~~#2%
}

% ── Titres ───────────────────────────────────────────────────────────────────
\usepackage{titlesec}
\titleformat{\section}  {\titlefont\bfseries\Huge\raggedright}  {}{0pt}{}
\titleformat{\subsection}{\titlefont\bfseries\Large\raggedright}{}{0pt}{}
\titlespacing{\section}  {0pt}{1.5em}{1em}
\titlespacing{\subsection}{0pt}{1em}{0.3em}

\newcommand{\seclabel}[2]{%
  {\small\titlefont\bfseries\color{#1}\MakeUppercase{#2}}\par\vspace{0.8em}%
}
\newcommand{\param}[1]{\textcolor{C2lo}{\texttt{#1}}}

% ── Fonds de page (eso-pic) ───────────────────────────────────────────────────
\usepackage{eso-pic}
\usepackage{pdflscape}

% Requiert : \documentclass déjà déclaré

% ── Police ───────────────────────────────────────────────────────────────────
\usepackage[T1]{fontenc}
\usepackage[utf8]{inputenc}
\usepackage{tgheros}
\renewcommand{\familydefault}{\sfdefault}
\newcommand{\titlefont}{\sffamily}

% ── Mise en page ──────────────────────────────────────────────────────────────
\usepackage[a4paper, top=2cm, bottom=2.2cm, left=2.5cm, right=2.5cm]{geometry}
\usepackage{microtype}
\setlength{\parindent}{0pt}
\setlength{\parskip}{0.6em}
\pagestyle{empty}

% ── Couleurs ─────────────────────────────────────────────────────────────────
\usepackage{xcolor}
\definecolor{C1lo}{HTML}{1E3A8A}  % bleu foncé
\definecolor{C1hi}{HTML}{0D9488}  % turquoise
\definecolor{C2lo}{HTML}{F4801A}  % orange
\definecolor{C2hi}{HTML}{C2185B}  % magenta
\definecolor{C3lo}{HTML}{27AE60}  % vert
\definecolor{C3hi}{HTML}{0D9488}  % turquoise (alias C1hi)
\definecolor{Cdark}{HTML}{1A1A2E}
\definecolor{Clight}{HTML}{F3F4F6}
\definecolor{Ccode}{HTML}{F0F0F0}

\definecolor{Cgrey}{HTML}{868686}
\colorlet{grey}{Cgrey}
\definecolor{Cc1}{HTML}{F0F0F0}
\definecolor{Cc2}{HTML}{F0F0F0}

% ── Mise en page avancée ─────────────────────────────────────────────────────
\usepackage{tabularx}
\usepackage{calc}
\usepackage{wrapfig}

% ── Commandes stat ────────────────────────────────────────────────────────────
%
% Six commandes de mise en page côte-à-côte : grand chiffre + sous-titre
% dans une colonne (gauche ou droite), texte courant dans l'autre colonne.
%
% Principe commun (variantes vtop/hbox)
% ──────────────────────────────────────
%   \hbox to \linewidth{ vtop(stat) vtop(texte) \hss }
%
%   vtop avec \kern0pt : point de référence au sommet (height=0, depth=contenu).
%   \parskip=0pt dans le vtop stat empêche un saut parasite avant le gros chiffre.
%
%   Espacement — variable \statgap
%   ──────────────────────────────
%   \statgap = blanc visible entre le chiffre et le texte courant.
%   Commun aux 6 commandes : changer \statgap recalibre tout.
%   Implémentation vtop/hbox : le blanc est absorbé dans la largeur du vtop-stat
%   (statwidth = natural + \statgap, pas de \hspace explicite entre les vtops).
%   Implémentation wrapfigure : \columnsep = \statgap (gap géré par wrapfig).
%
%   Alignement optique (sommet des caractères)
%   ───────────────────────────────────────────
%   Le vtop-texte est relevé par \raise\dimexpr\statcap-\bodycap\relax pour
%   aligner les sommets de glyphes. Non appliqué aux variantes wrapfigure.
%
% Comportement d'ancrage
% ──────────────────────
%   Les variantes vtop/hbox produisent un \hbox nu. Le caller gère \par et blancs.
%
% Les six commandes
% ─────────────────
%   \statcol    stat GAUCHE, vtop/hbox
%   \statflow   stat GAUCHE, vtop/hbox  (alias fonctionnel de \statcol)
%   \statflowr  stat DROITE, vtop/hbox, raggedleft, chiffre à la marge droite
%   \statcolwf  stat GAUCHE, wrapfigure (float : hors vbox/minipage)
%   \statflowwf stat GAUCHE, wrapfigure (idem, alias de \statcolwf)
%   \statflowrwf stat DROITE, wrapfigure (idem)

\newlength{\statwidth}
\newlength{\statcap}
\newlength{\bodycap}
\newlength{\statgap}
\setlength{\statgap}{1.0cm}

% \statcol{chiffre}{sous-titre}{texte}
\newcommand{\statcol}[3]{%
  \settowidth{\statwidth}{{\fontsize{60}{64}\selectfont\bfseries #1}}%
  \addtolength{\statwidth}{\statgap}%
  \settoheight{\statcap}{{\fontsize{60}{64}\selectfont\bfseries X}}%
  \settoheight{\bodycap}{{\large X}}%
  \noindent\hbox to \linewidth{%
    \vtop{\kern0pt\hsize=\statwidth\relax
      \parskip=0pt\relax
      {\fontsize{60}{64}\selectfont\bfseries #1\par}%
      {\large #2}%
    }%
    \raise\dimexpr\statcap-\bodycap\relax
    \vtop{\kern0pt\hsize=\dimexpr\linewidth-\statwidth\relax
      {\large #3\par}%
    }%
    \hss
  }%
}

% \statflow{chiffre}{sous-titre}{texte}  — stat gauche
\newcommand{\statflow}[3]{%
  \settowidth{\statwidth}{{\fontsize{60}{64}\selectfont\bfseries #1}}%
  \addtolength{\statwidth}{\statgap}%
  \settoheight{\statcap}{{\fontsize{60}{64}\selectfont\bfseries X}}%
  \settoheight{\bodycap}{{\large X}}%
  \noindent\hbox to \linewidth{%
    \vtop{\kern0pt\hsize=\statwidth\relax
      \parskip=0pt\relax
      {\fontsize{60}{64}\selectfont\bfseries #1\par}%
      {\large #2}%
    }%
    \raise\dimexpr\statcap-\bodycap\relax
    \vtop{\kern0pt\hsize=\dimexpr\linewidth-\statwidth\relax
      {\large #3\par}%
    }%
    \hss
  }%
}

% \statflowr{chiffre}{sous-titre}{texte}  — stat droite
\newcommand{\statflowr}[3]{%
  \settowidth{\statwidth}{{\fontsize{60}{64}\selectfont\bfseries #1}}%
  \addtolength{\statwidth}{\statgap}%
  \settoheight{\statcap}{{\fontsize{60}{64}\selectfont\bfseries X}}%
  \settoheight{\bodycap}{{\large X}}%
  \noindent\hbox to \linewidth{%
    \raise\dimexpr\statcap-\bodycap\relax
    \vtop{\kern0pt\hsize=\dimexpr\linewidth-\statwidth\relax
      {\large #3\par}%
    }%
    \vtop{\kern0pt\hsize=\statwidth\relax
      \parskip=0pt\relax
      \raggedleft
      {\fontsize{60}{64}\selectfont\bfseries #1\par}%
      {\large #2\par}%
    }%
    \hss
  }%
}

% ── Variantes wrapfigure ──────────────────────────────────────────────────────
%
% Doublons de statcol/statflow/statflowr utilisant wrapfigure.
% Le texte (#3) s'écoule autour du chiffre ; le contenu qui suit la commande
% dans le document wrappera aussi (comportement natif wrapfigure).
%
% Contrainte : wrapfigure est un float. Il ne fonctionne pas dans un vbox,
% minipage ou toute boîte : le float est alors différé hors de la boîte.
% À réserver au flot normal du document.
%
% Gap : \columnsep mis à \statgap en local.
% \intextsep mis à 0pt pour ne pas creuser de vide au-dessus du chiffre.
% La largeur déclarée à wrapfigure = natural_width du chiffre seul (le gap
% est pris en charge par \columnsep, pas inclus dans la boîte figure).

% \statcolwf{chiffre}{sous-titre}{texte}  — stat gauche, wrapfigure
\newcommand{\statcolwf}[3]{%
  \settowidth{\statwidth}{{\fontsize{60}{64}\selectfont\bfseries #1}}%
  \setlength{\intextsep}{0pt}%
  \setlength{\columnsep}{\statgap}%
  \begin{wrapfigure}{l}{\statwidth}%
    \vtop{\kern0pt\hsize=\statwidth\relax
      \parskip=0pt\relax
      {\fontsize{60}{64}\selectfont\bfseries #1\par}%
      {\large #2}%
    }%
  \end{wrapfigure}%
  {\large #3\par}%
}

% \statflowwf{chiffre}{sous-titre}{texte}  — stat gauche, wrapfigure
\newcommand{\statflowwf}[3]{%
  \settowidth{\statwidth}{{\fontsize{60}{64}\selectfont\bfseries #1}}%
  \setlength{\intextsep}{0pt}%
  \setlength{\columnsep}{\statgap}%
  \begin{wrapfigure}{l}{\statwidth}%
    \vtop{\kern0pt\hsize=\statwidth\relax
      \parskip=0pt\relax
      {\fontsize{60}{64}\selectfont\bfseries #1\par}%
      {\large #2}%
    }%
  \end{wrapfigure}%
  {\large #3\par}%
}

% \statflowrwf{chiffre}{sous-titre}{texte}  — stat droite, wrapfigure
\newcommand{\statflowrwf}[3]{%
  \settowidth{\statwidth}{{\fontsize{60}{64}\selectfont\bfseries #1}}%
  \setlength{\intextsep}{0pt}%
  \setlength{\columnsep}{\statgap}%
  \begin{wrapfigure}{r}{\statwidth}%
    \vtop{\kern0pt\hsize=\statwidth\relax
      \parskip=0pt\relax
      \raggedleft
      {\fontsize{60}{64}\selectfont\bfseries #1\par}%
      {\large #2\par}%
    }%
  \end{wrapfigure}%
  {\large #3\par}%
}

% ── TikZ ─────────────────────────────────────────────────────────────────────
\usepackage{tikz}
\usetikzlibrary{shadings, shadows, fadings}

% \gradtext{couleur-gauche}{couleur-droite}{contenu}
% Texte en dégradé gauche → droite, sur n'importe quel fond.
% Méthode : PDF Tr=7 (clip aux contours des glyphes) + dégradé TikZ.
% Après Q, le tracker CTM de pdflatex est désynchronisé ; on réémet un cm
% compensatoire égal à la largeur de l'image TikZ (convertie en bp).
\makeatletter
\newcommand{\gradtext}[3]{%
  \leavevmode%
  \pdfliteral{q 7 Tr }%
  \rlap{{#3}}%
  \pdfliteral{0 Tr }%
  \setbox\z@=\hbox{%
    \begin{tikzpicture}[baseline=(gtnode.base)]%
      \node[inner sep=0pt](gtnode){\phantom{#3}};%
      \shade[left color=#1, right color=#2]%
        (gtnode.south west)rectangle(gtnode.north east);%
    \end{tikzpicture}%
  }%
  \dimen@=\dimexpr\wd\z@ * 7200 / 7227\relax%
  \edef\gt@picwd{\strip@pt\dimen@}%
  \box\z@%
  \pdfliteral{Q}%
  \pdfliteral{1 0 0 1 \gt@picwd\space 0 cm}%
}
\makeatother

% ── Dégradés nommés (utilisés par \diagbg et \layeredbg) ─────────────────────
% sweep4 : froid → chaud (bleu / turquoise / orange / magenta)
\pgfdeclarehorizontalshading{sweep4}{100bp}{%
  color(0bp)=(C1lo);
  color(33bp)=(C1hi);
  color(66bp)=(C2lo);
  color(100bp)=(C2hi)%
}
% sweep3 : bleu → vert → orange
\pgfdeclarehorizontalshading{sweep3}{100bp}{%
  color(0bp)=(C1lo);
  color(50bp)=(C3lo);
  color(100bp)=(C2lo)%
}
% nightdawn : sombre → bleu → turquoise → orange
\pgfdeclarehorizontalshading{nightdawn}{100bp}{%
  color(0bp)=(Cdark);
  color(30bp)=(C1lo);
  color(60bp)=(C1hi);
  color(100bp)=(C2lo)%
}
% dusk : magenta → bleu → turquoise
\pgfdeclarehorizontalshading{dusk}{100bp}{%
  color(0bp)=(C2hi);
  color(45bp)=(C1lo);
  color(100bp)=(C1hi)%
}

% ── Anneaux radiaux (outer = white : fond blanc obligatoire) ──────────────────
\pgfdeclareradialshading{ringsA}{\pgfpointorigin}{%
  color(0.0cm)=(white!20!C1hi);
  color(0.9cm)=(C1hi);
  color(2.0cm)=(C1lo);
  color(3.4cm)=(C1lo!30!gray);
  color(4.8cm)=(gray!45);
  color(6.2cm)=(gray!18);
  color(7.6cm)=(gray!6);
  color(9.0cm)=(white)%
}
\pgfdeclareradialshading{ringsB}{\pgfpointorigin}{%
  color(0.0cm)=(white!20!C2hi);
  color(0.9cm)=(C2hi);
  color(2.0cm)=(C2lo);
  color(3.4cm)=(C2lo!30!gray);
  color(4.8cm)=(gray!45);
  color(6.2cm)=(gray!18);
  color(7.6cm)=(gray!6);
  color(9.0cm)=(white)%
}
\pgfdeclareradialshading{ringsC}{\pgfpointorigin}{%
  color(0.0cm)=(white!20!C3hi);
  color(0.9cm)=(C3hi);
  color(2.0cm)=(C3lo);
  color(3.4cm)=(C3lo!30!gray);
  color(4.8cm)=(gray!45);
  color(6.2cm)=(gray!18);
  color(7.6cm)=(gray!6);
  color(9.0cm)=(white)%
}

% ── Fonds de page ─────────────────────────────────────────────────────────────

% \placerings{ringsA|B|C}{anchor}{xcm}{ycm}
\newcommand{\placerings}[4]{%
  \AddToShipoutPictureBG*{%
    \begin{tikzpicture}[remember picture, overlay]
      \begin{scope}
        \clip (current page.south west) rectangle (current page.north east);
        \begin{scope}[shift={([xshift=#3 cm, yshift=#4 cm]current page.#2)}]
          \pgfuseshading{#1}
        \end{scope}
      \end{scope}
    \end{tikzpicture}%
  }%
}

% \spotfade{color}{xcm}{ycm}{anchor}{maxradius}
% À appeler à l'intérieur d'un tikzpicture remember picture/overlay.
\newcommand{\spotfade}[5]{%
  \begin{scope}
    \clip (current page.south west) rectangle (current page.north east);
    \pgfmathsetmacro{\Rmax}{#5}
    \foreach \frac/\op in {%
        1.00/0.04, 0.88/0.09, 0.76/0.14,
        0.64/0.20, 0.52/0.28, 0.40/0.38,
        0.28/0.49, 0.18/0.61, 0.10/0.72, 0.04/0.82}{%
      \pgfmathsetmacro{\r}{\Rmax*\frac}
      \fill[#1, opacity=\op]
        ([xshift=#2 cm, yshift=#3 cm]current page.#4) circle(\r cm);
    }
  \end{scope}
}

% \complexbg{baselo}{basehi}{spotcolor}{xcm}{ycm}{anchor}{radius}
\newcommand{\complexbg}[7]{%
  \AddToShipoutPictureBG*{%
    \begin{tikzpicture}[remember picture, overlay]
      \shade[top color=#1, bottom color=#2]
        (current page.south west) rectangle (current page.north east);
      \spotfade{#3}{#4}{#5}{#6}{#7}
    \end{tikzpicture}%
  }%
}

% \fourcornerbg{SW}{NW}{SE}{NE}
\newcommand{\fourcornerbg}[4]{%
  \AddToShipoutPictureBG*{%
    \begin{tikzpicture}[remember picture, overlay]
      \shade[lower left=#1, upper left=#2, lower right=#3, upper right=#4]
        (current page.south west) rectangle (current page.north east);
    \end{tikzpicture}%
  }%
}

% \halftopbg{SW}{NW}{SE}{NE}  — moitié haute colorée (bilinéaire), bas blanc
% Liseré solide en bas de page reprenant le dégradé gauche/droite du fond.
\newcommand{\halftopbg}[4]{%
  \AddToShipoutPictureBG*{%
    \begin{tikzpicture}[remember picture, overlay]
      \shade[lower left=#1, upper left=#2, lower right=#3, upper right=#4]
        (current page.west) rectangle (current page.north east);
      \shade[left color=#1, right color=#3]
        (current page.south west)
        rectangle
        ([yshift=12pt]current page.south east);
    \end{tikzpicture}%
  }%
}

% \topbg{splitht}{SW}{NW}{SE}{NE}  — haut coloré jusqu'à splitht depuis bas papier
% splitht : hauteur depuis current page.south (ex. 129.4mm pour alignement custom)
\newcommand{\topbg}[5]{%
  \AddToShipoutPictureBG*{%
    \begin{tikzpicture}[remember picture, overlay]
      \shade[lower left=#2, upper left=#3, lower right=#4, upper right=#5]
        ([yshift=#1]current page.south west)
        rectangle
        (current page.north east);
      \shade[left color=#2, right color=#4]
        (current page.south west)
        rectangle
        ([yshift=12pt]current page.south east);
    \end{tikzpicture}%
  }%
}

% \topbgsplit{splitht}{SW}{NW}{SE}{NE}  — comme \topbg, sans liseré bas (dessous blanc)
% splitht peut être \zposy{label}sp pour un split dynamique via zref-savepos.
\newcommand{\topbgsplit}[5]{%
  \AddToShipoutPictureBG*{%
    \begin{tikzpicture}[remember picture, overlay]
      \shade[lower left=#2, upper left=#3, lower right=#4, upper right=#5]
        ([yshift=#1]current page.south west)
        rectangle
        (current page.north east);
    \end{tikzpicture}%
  }%
}

% \halfbotbg{SW}{NW}{SE}{NE}  — moitié basse colorée, haut blanc
\newcommand{\halfbotbg}[4]{%
  \AddToShipoutPictureBG*{%
    \begin{tikzpicture}[remember picture, overlay]
      \shade[lower left=#1, upper left=#2, lower right=#3, upper right=#4]
        (current page.south west) rectangle (current page.east);
    \end{tikzpicture}%
  }%
}

% \diagbg{sweep4|sweep3|nightdawn|dusk}{angle}
%   angle : 0 (→)  45 (↗)  90 (↑)  135 (↖)  180 (←)
\newcommand{\diagbg}[2]{%
  \AddToShipoutPictureBG*{%
    \begin{tikzpicture}[remember picture, overlay]
      \shade[shading=#1, shading angle=#2]
        (current page.south west) rectangle (current page.north east);
    \end{tikzpicture}%
  }%
}

% \layeredbg{shading1}{angle1}{shading2}{angle2}{opacity2}
% Couche 2 à opacity2 par-dessus couche 1 : angles croisés → 4 coins distincts.
\newcommand{\layeredbg}[5]{%
  \AddToShipoutPictureBG*{%
    \begin{tikzpicture}[remember picture, overlay]
      \shade[shading=#1, shading angle=#2]
        (current page.south west) rectangle (current page.north east);
      \begin{scope}[opacity=#5]
        \shade[shading=#3, shading angle=#4]
          (current page.south west) rectangle (current page.north east);
      \end{scope}
    \end{tikzpicture}%
  }%
}

% \fourcornerspot{SW}{NW}{SE}{NE}{spotcolor}{radius}
\newcommand{\fourcornerspot}[6]{%
  \AddToShipoutPictureBG*{%
    \begin{tikzpicture}[remember picture, overlay]
      \shade[lower left=#1, upper left=#2, lower right=#3, upper right=#4]
        (current page.south west) rectangle (current page.north east);
      \spotfade{#5}{0}{0}{center}{#6}
    \end{tikzpicture}%
  }%
}

% \trispot{c-NW}{c-NE}{c-S}   fond sombre + halos NW / NE / bas-centre
\newcommand{\trispot}[3]{%
  \AddToShipoutPictureBG*{%
    \begin{tikzpicture}[remember picture, overlay]
      \fill[Cdark] (current page.south west) rectangle (current page.north east);
      \spotfade{#1}{0}{0}{north west}{12}
      \spotfade{#2}{0}{0}{north east}{12}
      \spotfade{#3}{0}{0}{south}{10}
    \end{tikzpicture}%
  }%
}

% \tricorner{c-NW}{c-NE}{c-SE}   fond sombre + halos NW / NE / SE
\newcommand{\tricorner}[3]{%
  \AddToShipoutPictureBG*{%
    \begin{tikzpicture}[remember picture, overlay]
      \fill[Cdark] (current page.south west) rectangle (current page.north east);
      \spotfade{#1}{0}{0}{north west}{12}
      \spotfade{#2}{0}{0}{north east}{12}
      \spotfade{#3}{0}{0}{south east}{12}
    \end{tikzpicture}%
  }%
}

% \tricornergb{c-NW}{c-NE}{c-SE}   gradient bilinéaire 3 coins, SW fixé à Cdark
\newcommand{\tricornergb}[3]{\fourcornerbg{Cdark}{#1}{#3}{#2}}

% \trispotgb{c-NW}{c-NE}{c-bas}   gradient bilinéaire, coins bas identiques
\newcommand{\trispotgb}[3]{\fourcornerbg{#3}{#1}{#3}{#2}}

% LaTeX limite à 9 args/commande : \doublebgA stocke les params, \doublebgB déclenche le rendu.
% \doublebgA{baselo}{basehi}{spot1color}{s1x}{s1y}{s1anchor}{s1radius}
% \doublebgB{spot2color}{s2x}{s2y}{s2anchor}{s2radius}
\newcommand{\doublebgA}[7]{%
  \gdef\dbblo{#1}\gdef\dbbhi{#2}%
  \gdef\dbca{#3}\gdef\dbxa{#4}\gdef\dbya{#5}\gdef\dbancha{#6}\gdef\dbra{#7}%
}
\newcommand{\doublebgB}[5]{%
  \AddToShipoutPictureBG*{%
    \begin{tikzpicture}[remember picture, overlay]
      \shade[top color=\dbblo, bottom color=\dbbhi]
        (current page.south west) rectangle (current page.north east);
      \spotfade{\dbca}{\dbxa}{\dbya}{\dbancha}{\dbra}
      \spotfade{#1}{#2}{#3}{#4}{#5}
    \end{tikzpicture}%
  }%
}

% ── Boites tcolorbox ─────────────────────────────────────────────────────────
\usepackage[most]{tcolorbox}
\tcbuselibrary{skins, breakable}

\newtcolorbox{gradbox}[2]{
  enhanced, breakable, arc=0pt, boxrule=0pt,
  interior style={left color=#1, right color=#2}, coltext=white,
  left=1.5cm, right=1.5cm, top=1cm, bottom=1cm,
  before skip=0.9em, after skip=0.9em,
}
\newtcolorbox{fullband}[2]{
  enhanced, breakable, arc=0pt, boxrule=0pt,
  interior style={left color=#1, right color=#2}, coltext=white,
  enlarge left by=-2.5cm, enlarge right by=-2.5cm,
  left=4.1cm, right=4.1cm, top=1cm, bottom=1cm,
  before skip=0.9em, after skip=0.9em,
}
\newtcolorbox{darkbox}{
  enhanced, breakable, arc=0pt, boxrule=0pt,
  colback=Cdark, colframe=Cdark, coltext=white,
  left=1.5cm, right=1.5cm, top=1cm, bottom=1cm,
  before skip=0.9em, after skip=0.9em,
}
\newtcolorbox{lightbox}{
  enhanced, breakable, arc=3pt, boxrule=0pt,
  colback=Clight, colframe=Clight,
  left=1.2cm, right=1.2cm, top=0.8cm, bottom=0.8cm,
  before skip=0.9em, after skip=0.9em,
}
\newtcolorbox{solidbox}[1]{
  enhanced, breakable, arc=0pt, boxrule=0pt,
  colback=#1, colframe=#1, coltext=white,
  left=1cm, right=1cm, top=0.8cm, bottom=0.8cm,
}
\newtcolorbox{whitebox}{
  enhanced, breakable, arc=3pt, boxrule=0pt,
  colback=white, colframe=white, coltext=black,
  drop shadow,
  left=1.2cm, right=1.2cm, top=0.9cm, bottom=0.9cm,
  before skip=0.9em, after skip=0.9em,
}
\newtcolorbox{codebox}{
  enhanced, breakable, arc=3pt, boxrule=0.4pt,
  colback=Ccode, colframe=gray!35, coltext=black,
  left=0.9cm, right=0.9cm, top=0.6cm, bottom=0.6cm,
  fontupper=\small\ttfamily,
  before skip=0.9em, after skip=0.9em,
}

\newcommand{\nbullet}[2]{%
  \begin{tikzpicture}[baseline=(c.base)]
    \node[circle, minimum size=1.6em, inner sep=0pt,
          top color=C1lo, bottom color=C1hi,
          text=white, font=\bfseries\small] (c) {#1};
  \end{tikzpicture}~~#2%
}

% ── Titres ───────────────────────────────────────────────────────────────────
\usepackage{titlesec}
\titleformat{\section}  {\titlefont\bfseries\Huge\raggedright}  {}{0pt}{}
\titleformat{\subsection}{\titlefont\bfseries\Large\raggedright}{}{0pt}{}
\titlespacing{\section}  {0pt}{1.5em}{1em}
\titlespacing{\subsection}{0pt}{1em}{0.3em}

\newcommand{\seclabel}[2]{%
  {\small\titlefont\bfseries\color{#1}\MakeUppercase{#2}}\par\vspace{0.8em}%
}
\newcommand{\param}[1]{\textcolor{C2lo}{\texttt{#1}}}

% ── Fonds de page (eso-pic) ───────────────────────────────────────────────────
\usepackage{eso-pic}
\usepackage{pdflscape}

% Requiert : \documentclass déjà déclaré

% ── Police ───────────────────────────────────────────────────────────────────
\usepackage[T1]{fontenc}
\usepackage[utf8]{inputenc}
\usepackage{tgheros}
\renewcommand{\familydefault}{\sfdefault}
\newcommand{\titlefont}{\sffamily}

% ── Mise en page ──────────────────────────────────────────────────────────────
\usepackage[a4paper, top=2cm, bottom=2.2cm, left=2.5cm, right=2.5cm]{geometry}
\usepackage{microtype}
\setlength{\parindent}{0pt}
\setlength{\parskip}{0.6em}
\pagestyle{empty}

% ── Couleurs ─────────────────────────────────────────────────────────────────
\usepackage{xcolor}
\definecolor{C1lo}{HTML}{1E3A8A}  % bleu foncé
\definecolor{C1hi}{HTML}{0D9488}  % turquoise
\definecolor{C2lo}{HTML}{F4801A}  % orange
\definecolor{C2hi}{HTML}{C2185B}  % magenta
\definecolor{C3lo}{HTML}{27AE60}  % vert
\definecolor{C3hi}{HTML}{0D9488}  % turquoise (alias C1hi)
\definecolor{Cdark}{HTML}{1A1A2E}
\definecolor{Clight}{HTML}{F3F4F6}
\definecolor{Ccode}{HTML}{F0F0F0}

\definecolor{Cgrey}{HTML}{868686}
\colorlet{grey}{Cgrey}
\definecolor{Cc1}{HTML}{F0F0F0}
\definecolor{Cc2}{HTML}{F0F0F0}

% ── Mise en page avancée ─────────────────────────────────────────────────────
\usepackage{tabularx}
\usepackage{calc}
\usepackage{wrapfig}

% ── Commandes stat ────────────────────────────────────────────────────────────
%
% Six commandes de mise en page côte-à-côte : grand chiffre + sous-titre
% dans une colonne (gauche ou droite), texte courant dans l'autre colonne.
%
% Principe commun (variantes vtop/hbox)
% ──────────────────────────────────────
%   \hbox to \linewidth{ vtop(stat) vtop(texte) \hss }
%
%   vtop avec \kern0pt : point de référence au sommet (height=0, depth=contenu).
%   \parskip=0pt dans le vtop stat empêche un saut parasite avant le gros chiffre.
%
%   Espacement — variable \statgap
%   ──────────────────────────────
%   \statgap = blanc visible entre le chiffre et le texte courant.
%   Commun aux 6 commandes : changer \statgap recalibre tout.
%   Implémentation vtop/hbox : le blanc est absorbé dans la largeur du vtop-stat
%   (statwidth = natural + \statgap, pas de \hspace explicite entre les vtops).
%   Implémentation wrapfigure : \columnsep = \statgap (gap géré par wrapfig).
%
%   Alignement optique (sommet des caractères)
%   ───────────────────────────────────────────
%   Le vtop-texte est relevé par \raise\dimexpr\statcap-\bodycap\relax pour
%   aligner les sommets de glyphes. Non appliqué aux variantes wrapfigure.
%
% Comportement d'ancrage
% ──────────────────────
%   Les variantes vtop/hbox produisent un \hbox nu. Le caller gère \par et blancs.
%
% Les six commandes
% ─────────────────
%   \statcol    stat GAUCHE, vtop/hbox
%   \statflow   stat GAUCHE, vtop/hbox  (alias fonctionnel de \statcol)
%   \statflowr  stat DROITE, vtop/hbox, raggedleft, chiffre à la marge droite
%   \statcolwf  stat GAUCHE, wrapfigure (float : hors vbox/minipage)
%   \statflowwf stat GAUCHE, wrapfigure (idem, alias de \statcolwf)
%   \statflowrwf stat DROITE, wrapfigure (idem)

\newlength{\statwidth}
\newlength{\statcap}
\newlength{\bodycap}
\newlength{\statgap}
\setlength{\statgap}{1.0cm}

% \statcol{chiffre}{sous-titre}{texte}
\newcommand{\statcol}[3]{%
  \settowidth{\statwidth}{{\fontsize{60}{64}\selectfont\bfseries #1}}%
  \addtolength{\statwidth}{\statgap}%
  \settoheight{\statcap}{{\fontsize{60}{64}\selectfont\bfseries X}}%
  \settoheight{\bodycap}{{\large X}}%
  \noindent\hbox to \linewidth{%
    \vtop{\kern0pt\hsize=\statwidth\relax
      \parskip=0pt\relax
      {\fontsize{60}{64}\selectfont\bfseries #1\par}%
      {\large #2}%
    }%
    \raise\dimexpr\statcap-\bodycap\relax
    \vtop{\kern0pt\hsize=\dimexpr\linewidth-\statwidth\relax
      {\large #3\par}%
    }%
    \hss
  }%
}

% \statflow{chiffre}{sous-titre}{texte}  — stat gauche
\newcommand{\statflow}[3]{%
  \settowidth{\statwidth}{{\fontsize{60}{64}\selectfont\bfseries #1}}%
  \addtolength{\statwidth}{\statgap}%
  \settoheight{\statcap}{{\fontsize{60}{64}\selectfont\bfseries X}}%
  \settoheight{\bodycap}{{\large X}}%
  \noindent\hbox to \linewidth{%
    \vtop{\kern0pt\hsize=\statwidth\relax
      \parskip=0pt\relax
      {\fontsize{60}{64}\selectfont\bfseries #1\par}%
      {\large #2}%
    }%
    \raise\dimexpr\statcap-\bodycap\relax
    \vtop{\kern0pt\hsize=\dimexpr\linewidth-\statwidth\relax
      {\large #3\par}%
    }%
    \hss
  }%
}

% \statflowr{chiffre}{sous-titre}{texte}  — stat droite
\newcommand{\statflowr}[3]{%
  \settowidth{\statwidth}{{\fontsize{60}{64}\selectfont\bfseries #1}}%
  \addtolength{\statwidth}{\statgap}%
  \settoheight{\statcap}{{\fontsize{60}{64}\selectfont\bfseries X}}%
  \settoheight{\bodycap}{{\large X}}%
  \noindent\hbox to \linewidth{%
    \raise\dimexpr\statcap-\bodycap\relax
    \vtop{\kern0pt\hsize=\dimexpr\linewidth-\statwidth\relax
      {\large #3\par}%
    }%
    \vtop{\kern0pt\hsize=\statwidth\relax
      \parskip=0pt\relax
      \raggedleft
      {\fontsize{60}{64}\selectfont\bfseries #1\par}%
      {\large #2\par}%
    }%
    \hss
  }%
}

% ── Variantes wrapfigure ──────────────────────────────────────────────────────
%
% Doublons de statcol/statflow/statflowr utilisant wrapfigure.
% Le texte (#3) s'écoule autour du chiffre ; le contenu qui suit la commande
% dans le document wrappera aussi (comportement natif wrapfigure).
%
% Contrainte : wrapfigure est un float. Il ne fonctionne pas dans un vbox,
% minipage ou toute boîte : le float est alors différé hors de la boîte.
% À réserver au flot normal du document.
%
% Gap : \columnsep mis à \statgap en local.
% \intextsep mis à 0pt pour ne pas creuser de vide au-dessus du chiffre.
% La largeur déclarée à wrapfigure = natural_width du chiffre seul (le gap
% est pris en charge par \columnsep, pas inclus dans la boîte figure).

% \statcolwf{chiffre}{sous-titre}{texte}  — stat gauche, wrapfigure
\newcommand{\statcolwf}[3]{%
  \settowidth{\statwidth}{{\fontsize{60}{64}\selectfont\bfseries #1}}%
  \setlength{\intextsep}{0pt}%
  \setlength{\columnsep}{\statgap}%
  \begin{wrapfigure}{l}{\statwidth}%
    \vtop{\kern0pt\hsize=\statwidth\relax
      \parskip=0pt\relax
      {\fontsize{60}{64}\selectfont\bfseries #1\par}%
      {\large #2}%
    }%
  \end{wrapfigure}%
  {\large #3\par}%
}

% \statflowwf{chiffre}{sous-titre}{texte}  — stat gauche, wrapfigure
\newcommand{\statflowwf}[3]{%
  \settowidth{\statwidth}{{\fontsize{60}{64}\selectfont\bfseries #1}}%
  \setlength{\intextsep}{0pt}%
  \setlength{\columnsep}{\statgap}%
  \begin{wrapfigure}{l}{\statwidth}%
    \vtop{\kern0pt\hsize=\statwidth\relax
      \parskip=0pt\relax
      {\fontsize{60}{64}\selectfont\bfseries #1\par}%
      {\large #2}%
    }%
  \end{wrapfigure}%
  {\large #3\par}%
}

% \statflowrwf{chiffre}{sous-titre}{texte}  — stat droite, wrapfigure
\newcommand{\statflowrwf}[3]{%
  \settowidth{\statwidth}{{\fontsize{60}{64}\selectfont\bfseries #1}}%
  \setlength{\intextsep}{0pt}%
  \setlength{\columnsep}{\statgap}%
  \begin{wrapfigure}{r}{\statwidth}%
    \vtop{\kern0pt\hsize=\statwidth\relax
      \parskip=0pt\relax
      \raggedleft
      {\fontsize{60}{64}\selectfont\bfseries #1\par}%
      {\large #2\par}%
    }%
  \end{wrapfigure}%
  {\large #3\par}%
}

% ── TikZ ─────────────────────────────────────────────────────────────────────
\usepackage{tikz}
\usetikzlibrary{shadings, shadows, fadings}

% \gradtext{couleur-gauche}{couleur-droite}{contenu}
% Texte en dégradé gauche → droite, sur n'importe quel fond.
% Méthode : PDF Tr=7 (clip aux contours des glyphes) + dégradé TikZ.
% Après Q, le tracker CTM de pdflatex est désynchronisé ; on réémet un cm
% compensatoire égal à la largeur de l'image TikZ (convertie en bp).
\makeatletter
\newcommand{\gradtext}[3]{%
  \leavevmode%
  \pdfliteral{q 7 Tr }%
  \rlap{{#3}}%
  \pdfliteral{0 Tr }%
  \setbox\z@=\hbox{%
    \begin{tikzpicture}[baseline=(gtnode.base)]%
      \node[inner sep=0pt](gtnode){\phantom{#3}};%
      \shade[left color=#1, right color=#2]%
        (gtnode.south west)rectangle(gtnode.north east);%
    \end{tikzpicture}%
  }%
  \dimen@=\dimexpr\wd\z@ * 7200 / 7227\relax%
  \edef\gt@picwd{\strip@pt\dimen@}%
  \box\z@%
  \pdfliteral{Q}%
  \pdfliteral{1 0 0 1 \gt@picwd\space 0 cm}%
}
\makeatother

% ── Dégradés nommés (utilisés par \diagbg et \layeredbg) ─────────────────────
% sweep4 : froid → chaud (bleu / turquoise / orange / magenta)
\pgfdeclarehorizontalshading{sweep4}{100bp}{%
  color(0bp)=(C1lo);
  color(33bp)=(C1hi);
  color(66bp)=(C2lo);
  color(100bp)=(C2hi)%
}
% sweep3 : bleu → vert → orange
\pgfdeclarehorizontalshading{sweep3}{100bp}{%
  color(0bp)=(C1lo);
  color(50bp)=(C3lo);
  color(100bp)=(C2lo)%
}
% nightdawn : sombre → bleu → turquoise → orange
\pgfdeclarehorizontalshading{nightdawn}{100bp}{%
  color(0bp)=(Cdark);
  color(30bp)=(C1lo);
  color(60bp)=(C1hi);
  color(100bp)=(C2lo)%
}
% dusk : magenta → bleu → turquoise
\pgfdeclarehorizontalshading{dusk}{100bp}{%
  color(0bp)=(C2hi);
  color(45bp)=(C1lo);
  color(100bp)=(C1hi)%
}

% ── Anneaux radiaux (outer = white : fond blanc obligatoire) ──────────────────
\pgfdeclareradialshading{ringsA}{\pgfpointorigin}{%
  color(0.0cm)=(white!20!C1hi);
  color(0.9cm)=(C1hi);
  color(2.0cm)=(C1lo);
  color(3.4cm)=(C1lo!30!gray);
  color(4.8cm)=(gray!45);
  color(6.2cm)=(gray!18);
  color(7.6cm)=(gray!6);
  color(9.0cm)=(white)%
}
\pgfdeclareradialshading{ringsB}{\pgfpointorigin}{%
  color(0.0cm)=(white!20!C2hi);
  color(0.9cm)=(C2hi);
  color(2.0cm)=(C2lo);
  color(3.4cm)=(C2lo!30!gray);
  color(4.8cm)=(gray!45);
  color(6.2cm)=(gray!18);
  color(7.6cm)=(gray!6);
  color(9.0cm)=(white)%
}
\pgfdeclareradialshading{ringsC}{\pgfpointorigin}{%
  color(0.0cm)=(white!20!C3hi);
  color(0.9cm)=(C3hi);
  color(2.0cm)=(C3lo);
  color(3.4cm)=(C3lo!30!gray);
  color(4.8cm)=(gray!45);
  color(6.2cm)=(gray!18);
  color(7.6cm)=(gray!6);
  color(9.0cm)=(white)%
}

% ── Fonds de page ─────────────────────────────────────────────────────────────

% \placerings{ringsA|B|C}{anchor}{xcm}{ycm}
\newcommand{\placerings}[4]{%
  \AddToShipoutPictureBG*{%
    \begin{tikzpicture}[remember picture, overlay]
      \begin{scope}
        \clip (current page.south west) rectangle (current page.north east);
        \begin{scope}[shift={([xshift=#3 cm, yshift=#4 cm]current page.#2)}]
          \pgfuseshading{#1}
        \end{scope}
      \end{scope}
    \end{tikzpicture}%
  }%
}

% \spotfade{color}{xcm}{ycm}{anchor}{maxradius}
% À appeler à l'intérieur d'un tikzpicture remember picture/overlay.
\newcommand{\spotfade}[5]{%
  \begin{scope}
    \clip (current page.south west) rectangle (current page.north east);
    \pgfmathsetmacro{\Rmax}{#5}
    \foreach \frac/\op in {%
        1.00/0.04, 0.88/0.09, 0.76/0.14,
        0.64/0.20, 0.52/0.28, 0.40/0.38,
        0.28/0.49, 0.18/0.61, 0.10/0.72, 0.04/0.82}{%
      \pgfmathsetmacro{\r}{\Rmax*\frac}
      \fill[#1, opacity=\op]
        ([xshift=#2 cm, yshift=#3 cm]current page.#4) circle(\r cm);
    }
  \end{scope}
}

% \complexbg{baselo}{basehi}{spotcolor}{xcm}{ycm}{anchor}{radius}
\newcommand{\complexbg}[7]{%
  \AddToShipoutPictureBG*{%
    \begin{tikzpicture}[remember picture, overlay]
      \shade[top color=#1, bottom color=#2]
        (current page.south west) rectangle (current page.north east);
      \spotfade{#3}{#4}{#5}{#6}{#7}
    \end{tikzpicture}%
  }%
}

% \fourcornerbg{SW}{NW}{SE}{NE}
\newcommand{\fourcornerbg}[4]{%
  \AddToShipoutPictureBG*{%
    \begin{tikzpicture}[remember picture, overlay]
      \shade[lower left=#1, upper left=#2, lower right=#3, upper right=#4]
        (current page.south west) rectangle (current page.north east);
    \end{tikzpicture}%
  }%
}

% \halftopbg{SW}{NW}{SE}{NE}  — moitié haute colorée (bilinéaire), bas blanc
% Liseré solide en bas de page reprenant le dégradé gauche/droite du fond.
\newcommand{\halftopbg}[4]{%
  \AddToShipoutPictureBG*{%
    \begin{tikzpicture}[remember picture, overlay]
      \shade[lower left=#1, upper left=#2, lower right=#3, upper right=#4]
        (current page.west) rectangle (current page.north east);
      \shade[left color=#1, right color=#3]
        (current page.south west)
        rectangle
        ([yshift=12pt]current page.south east);
    \end{tikzpicture}%
  }%
}

% \topbg{splitht}{SW}{NW}{SE}{NE}  — haut coloré jusqu'à splitht depuis bas papier
% splitht : hauteur depuis current page.south (ex. 129.4mm pour alignement custom)
\newcommand{\topbg}[5]{%
  \AddToShipoutPictureBG*{%
    \begin{tikzpicture}[remember picture, overlay]
      \shade[lower left=#2, upper left=#3, lower right=#4, upper right=#5]
        ([yshift=#1]current page.south west)
        rectangle
        (current page.north east);
      \shade[left color=#2, right color=#4]
        (current page.south west)
        rectangle
        ([yshift=12pt]current page.south east);
    \end{tikzpicture}%
  }%
}

% \topbgsplit{splitht}{SW}{NW}{SE}{NE}  — comme \topbg, sans liseré bas (dessous blanc)
% splitht peut être \zposy{label}sp pour un split dynamique via zref-savepos.
\newcommand{\topbgsplit}[5]{%
  \AddToShipoutPictureBG*{%
    \begin{tikzpicture}[remember picture, overlay]
      \shade[lower left=#2, upper left=#3, lower right=#4, upper right=#5]
        ([yshift=#1]current page.south west)
        rectangle
        (current page.north east);
    \end{tikzpicture}%
  }%
}

% \halfbotbg{SW}{NW}{SE}{NE}  — moitié basse colorée, haut blanc
\newcommand{\halfbotbg}[4]{%
  \AddToShipoutPictureBG*{%
    \begin{tikzpicture}[remember picture, overlay]
      \shade[lower left=#1, upper left=#2, lower right=#3, upper right=#4]
        (current page.south west) rectangle (current page.east);
    \end{tikzpicture}%
  }%
}

% \diagbg{sweep4|sweep3|nightdawn|dusk}{angle}
%   angle : 0 (→)  45 (↗)  90 (↑)  135 (↖)  180 (←)
\newcommand{\diagbg}[2]{%
  \AddToShipoutPictureBG*{%
    \begin{tikzpicture}[remember picture, overlay]
      \shade[shading=#1, shading angle=#2]
        (current page.south west) rectangle (current page.north east);
    \end{tikzpicture}%
  }%
}

% \layeredbg{shading1}{angle1}{shading2}{angle2}{opacity2}
% Couche 2 à opacity2 par-dessus couche 1 : angles croisés → 4 coins distincts.
\newcommand{\layeredbg}[5]{%
  \AddToShipoutPictureBG*{%
    \begin{tikzpicture}[remember picture, overlay]
      \shade[shading=#1, shading angle=#2]
        (current page.south west) rectangle (current page.north east);
      \begin{scope}[opacity=#5]
        \shade[shading=#3, shading angle=#4]
          (current page.south west) rectangle (current page.north east);
      \end{scope}
    \end{tikzpicture}%
  }%
}

% \fourcornerspot{SW}{NW}{SE}{NE}{spotcolor}{radius}
\newcommand{\fourcornerspot}[6]{%
  \AddToShipoutPictureBG*{%
    \begin{tikzpicture}[remember picture, overlay]
      \shade[lower left=#1, upper left=#2, lower right=#3, upper right=#4]
        (current page.south west) rectangle (current page.north east);
      \spotfade{#5}{0}{0}{center}{#6}
    \end{tikzpicture}%
  }%
}

% \trispot{c-NW}{c-NE}{c-S}   fond sombre + halos NW / NE / bas-centre
\newcommand{\trispot}[3]{%
  \AddToShipoutPictureBG*{%
    \begin{tikzpicture}[remember picture, overlay]
      \fill[Cdark] (current page.south west) rectangle (current page.north east);
      \spotfade{#1}{0}{0}{north west}{12}
      \spotfade{#2}{0}{0}{north east}{12}
      \spotfade{#3}{0}{0}{south}{10}
    \end{tikzpicture}%
  }%
}

% \tricorner{c-NW}{c-NE}{c-SE}   fond sombre + halos NW / NE / SE
\newcommand{\tricorner}[3]{%
  \AddToShipoutPictureBG*{%
    \begin{tikzpicture}[remember picture, overlay]
      \fill[Cdark] (current page.south west) rectangle (current page.north east);
      \spotfade{#1}{0}{0}{north west}{12}
      \spotfade{#2}{0}{0}{north east}{12}
      \spotfade{#3}{0}{0}{south east}{12}
    \end{tikzpicture}%
  }%
}

% \tricornergb{c-NW}{c-NE}{c-SE}   gradient bilinéaire 3 coins, SW fixé à Cdark
\newcommand{\tricornergb}[3]{\fourcornerbg{Cdark}{#1}{#3}{#2}}

% \trispotgb{c-NW}{c-NE}{c-bas}   gradient bilinéaire, coins bas identiques
\newcommand{\trispotgb}[3]{\fourcornerbg{#3}{#1}{#3}{#2}}

% LaTeX limite à 9 args/commande : \doublebgA stocke les params, \doublebgB déclenche le rendu.
% \doublebgA{baselo}{basehi}{spot1color}{s1x}{s1y}{s1anchor}{s1radius}
% \doublebgB{spot2color}{s2x}{s2y}{s2anchor}{s2radius}
\newcommand{\doublebgA}[7]{%
  \gdef\dbblo{#1}\gdef\dbbhi{#2}%
  \gdef\dbca{#3}\gdef\dbxa{#4}\gdef\dbya{#5}\gdef\dbancha{#6}\gdef\dbra{#7}%
}
\newcommand{\doublebgB}[5]{%
  \AddToShipoutPictureBG*{%
    \begin{tikzpicture}[remember picture, overlay]
      \shade[top color=\dbblo, bottom color=\dbbhi]
        (current page.south west) rectangle (current page.north east);
      \spotfade{\dbca}{\dbxa}{\dbya}{\dbancha}{\dbra}
      \spotfade{#1}{#2}{#3}{#4}{#5}
    \end{tikzpicture}%
  }%
}

% ── Boites tcolorbox ─────────────────────────────────────────────────────────
\usepackage[most]{tcolorbox}
\tcbuselibrary{skins, breakable}

\newtcolorbox{gradbox}[2]{
  enhanced, breakable, arc=0pt, boxrule=0pt,
  interior style={left color=#1, right color=#2}, coltext=white,
  left=1.5cm, right=1.5cm, top=1cm, bottom=1cm,
  before skip=0.9em, after skip=0.9em,
}
\newtcolorbox{fullband}[2]{
  enhanced, breakable, arc=0pt, boxrule=0pt,
  interior style={left color=#1, right color=#2}, coltext=white,
  enlarge left by=-2.5cm, enlarge right by=-2.5cm,
  left=4.1cm, right=4.1cm, top=1cm, bottom=1cm,
  before skip=0.9em, after skip=0.9em,
}
\newtcolorbox{darkbox}{
  enhanced, breakable, arc=0pt, boxrule=0pt,
  colback=Cdark, colframe=Cdark, coltext=white,
  left=1.5cm, right=1.5cm, top=1cm, bottom=1cm,
  before skip=0.9em, after skip=0.9em,
}
\newtcolorbox{lightbox}{
  enhanced, breakable, arc=3pt, boxrule=0pt,
  colback=Clight, colframe=Clight,
  left=1.2cm, right=1.2cm, top=0.8cm, bottom=0.8cm,
  before skip=0.9em, after skip=0.9em,
}
\newtcolorbox{solidbox}[1]{
  enhanced, breakable, arc=0pt, boxrule=0pt,
  colback=#1, colframe=#1, coltext=white,
  left=1cm, right=1cm, top=0.8cm, bottom=0.8cm,
}
\newtcolorbox{whitebox}{
  enhanced, breakable, arc=3pt, boxrule=0pt,
  colback=white, colframe=white, coltext=black,
  drop shadow,
  left=1.2cm, right=1.2cm, top=0.9cm, bottom=0.9cm,
  before skip=0.9em, after skip=0.9em,
}
\newtcolorbox{codebox}{
  enhanced, breakable, arc=3pt, boxrule=0.4pt,
  colback=Ccode, colframe=gray!35, coltext=black,
  left=0.9cm, right=0.9cm, top=0.6cm, bottom=0.6cm,
  fontupper=\small\ttfamily,
  before skip=0.9em, after skip=0.9em,
}

\newcommand{\nbullet}[2]{%
  \begin{tikzpicture}[baseline=(c.base)]
    \node[circle, minimum size=1.6em, inner sep=0pt,
          top color=C1lo, bottom color=C1hi,
          text=white, font=\bfseries\small] (c) {#1};
  \end{tikzpicture}~~#2%
}

% ── Titres ───────────────────────────────────────────────────────────────────
\usepackage{titlesec}
\titleformat{\section}  {\titlefont\bfseries\Huge\raggedright}  {}{0pt}{}
\titleformat{\subsection}{\titlefont\bfseries\Large\raggedright}{}{0pt}{}
\titlespacing{\section}  {0pt}{1.5em}{1em}
\titlespacing{\subsection}{0pt}{1em}{0.3em}

\newcommand{\seclabel}[2]{%
  {\small\titlefont\bfseries\color{#1}\MakeUppercase{#2}}\par\vspace{0.8em}%
}
\newcommand{\param}[1]{\textcolor{C2lo}{\texttt{#1}}}

% ── Fonds de page (eso-pic) ───────────────────────────────────────────────────
\usepackage{eso-pic}
\usepackage{pdflscape}
